% ID: SEL-D0033-OA-E04
% ESTADO: cotejada. Segunda pasada visual y borrador compilado revisados.
% Fuente: Selectividad/matematicas_aplicadas/titular_septiembre/Examen-Matemáticas Aplicadas a las Ciencias Sociales.pdf
% SHA-256: 0f95b6aefa78f1a13e2115b281338bd807177671c303c380180957fe03f0a605
% Página: 1; opción A; ejercicio original 4.

En una zona escolar formada por tres centros de secundaria, se desea estimar la proporción del alumnado que lleva teléfono móvil al instituto. Se toma una muestra aleatoria simple de 121 estudiantes, de los cuales 74 lo llevan.
\begin{enumerate}[label=\alph*)]
\item \textbf{(1.2 puntos)} Determine un intervalo de confianza al 97\% para la proporción de este alumnado que lleva el móvil al instituto. \textquestiondown{}Entre qué dos porcentajes varía esa proporción a ese nivel de confianza?
\item \textbf{(0.5 puntos)} Si con la misma muestra se disminuye el nivel de confianza, \textquestiondown{}qué efecto tendrá esta disminución en el error de estimación?
\item \textbf{(0.8 puntos)} Si en la misma zona se elige mediante muestreo aleatorio estratificado con afijación proporcional otra muestra de 121 estudiantes, considerando que el segundo centro escolar tiene el doble de alumnos que el primero y el tercero tiene el triple que el primero, \textquestiondown{}cuántos alumnos de cada centro se deben tomar para constituir la muestra?
\end{enumerate}
